\documentclass[a4paper]{article}
\usepackage[pdftex]{graphicx}
\usepackage[utf8]{inputenc}
\usepackage{enumerate}
\usepackage{siunitx}
\sisetup{locale=DE} 
\usepackage{href-ul}
\hypersetup{
	colorlinks=true,
	linkcolor=blue,
	urlcolor=blue}
\usepackage{icomma}
\usepackage{amssymb}
\usepackage{geometry}
\geometry{a4paper, top=15mm, left=15mm, right=15mm, bottom=15mm,
	headsep=10mm, footskip=12mm}
%Source: img_0234

% Start the document
\begin{document}
%Titel
{\bf Kondensator im Wechselstromkreis }\\
\begin{enumerate}[1.]

%Aufgabe 1
\begin{minipage}{0.65\textwidth}
\item An einer Reihenschaltung aus einem Kondensator mit der Kapazität C und einem ohmschen Widerstand $R=\SI{40}{\kilo\ohm}$ liegt die Eingangsspannung \\
$U(t)=\hat{U}\cdot \sin{(2 \pi f t)}$. Der Effektivwert $U_{eff}$ der Eingangsspannung $U$ ist konstant und beträgt \\ $U_{eff}=\SI{12}{\volt}$. Die Frequenz dieser Eingangsspannung $U$ ist veränderbar.
Als Ausgangsspannung kann man die Spannung $U_R$ am ohmschen Widerstand $R$ oder die Spannung am Kondensator abgreifen. Die Effektivwerte $U_{R,eff}$ und $U_{C,eff}$ sind abhängig von der Frequenz $f$ der Eingangsspannung $U$.
\end{minipage}
\hfill
\begin{minipage}{0.35\textwidth}
\includegraphics[width=5.5 cm]{weko234}
\end{minipage}

Zeigen Sie mit Hilfe eines Zeigerdiagramms für Spannungen, dass für den Scheinwiderstand $Z$ gilt: $$ Z= \sqrt{R^2 + {1 \over 4 \pi^2 C^2} \cdot {1 \over f^2}}$$

%Aufgabe 2
\item Die Frequenz $f$ der Eingangsspannung $U$ wird zunächst auf den Wert $f_1=\SI{520}{\hertz}$ eingestellt. Bei dieser Frequenz $f_1$ fließt im Wechselstromkreis ein Strom mit der Stärke $I_1$ und die am ohmschen Widerstand $R$ abfallende Spannung $U_{R,1}$ hat den Effektivwert 
$U_{R,eff,1}=\SI{11,2}{\volt}$ 
\begin{enumerate}[a)]
\item Berechnen Sie die Kapazität $C$ des Kondensators und die zwischen der Stromstärke $I_1$ und der Eingangsspannung auftretende Phasenverschiebung $\varphi_1$. $[\textnormal{Teilergebnis:}\quad C=\SI{20}{\nano\farad}]$
\item Berechnen Sie die mittlere Wirkleistung $P_1$ des Wechselstromkreises.
\end{enumerate}

%Aufgabe 3
\item Die Frequenz $f$ der Eingangsspannung $U$ wird nun variiert. Dabei wird zunächst die Spannung $U_R$ am ohmschen Widerstand als Ausgangsspannung abgegriffen.
\begin{enumerate}[a)]
\item Bestätigen Sie durch allgemeine Rechnung, dass für den von der Frequenz $f$ abhängigen Effektivwert $U_{R,eff}$ der Spannung $U_R$ gilt: 
$$U_{R,eff}(f)={R \over \sqrt{R^2 + {1 \over 4 \pi^2 C^2} \cdot {1 \over f^2}}}\cdot U_{eff}$$
\item Wird die Frequenz $f$ der Eingangsspannung $U$ auf den Wert $f_g$ eingestellt, so sind die Effektivwerte $U_{R,eff}$ und $U_{C,eff}$ gleich groß. Berechnen sie die so genannte Grenzfrequenz $f_g$ und die Effektivwerte $U_{R,eff}$ und $U_{C,eff}$ für diese Frequenz $f_g$. $[\textnormal{Teilergebnis:}\quad f_g=\SI{0,20}{\kilo\hertz}]$
\item Begründen Sie mit Hilfe der Formel von Teilaufgabe 3. a), dass für hohe Frequenzen 
$f (f \gg f_g)$ der Effektivwert  $U_{R,eff}$ der Ausgangsspannung $U_R$ praktisch so groß ist wie der Effektivwert $U_{eff}$ der Eingangsspannung $U$.
\end{enumerate}
\end{enumerate} 

\fbox{
	\begin{minipage}{0.5\textwidth}
		Zur Lösung bitte \href{https://www.okuyakl.de/phys/p13wekoL234/ll234.pdf}{hier klicken} oder den QR-Code scannen.\\
	Weitere Arbeitsblätter gibt es unter 
	
	\href{https://www.okuyakl.de}{www.okuyakl.de}
	\end{minipage}
	\hfill
	\begin{minipage}{0.4\textwidth}
		\includegraphics[width=1.5 cm]{../../viecher/zwe03}
		\includegraphics[width=3 cm]{qr234}
		\includegraphics[width=2 cm]{../../viecher/afanticon1}
		
	\end{minipage}}

\end{document}
